空气中一片理想薄对称双凸透镜的两面曲率半径分别为 $R_1=R>0$、$R_2=-R$。折射率随真空波长按正常色散规律 $n(\lambda)=n_\infty+C/\lambda^2$ 变化，其中 $n_\infty>1,C>0$。忽略厚度、像差、反射、吸收及高阶色散。点物体位于透镜左侧距离 $s>0$，并在参考波长 $\lambda_0>0$ 满足 $s>f(\lambda_0)$；像距 $s'(\lambda)$ 以右侧为正。
(1) 从透镜制造公式求焦距 $f(\lambda)$ 及其随波长变化的方向；(2) 用高斯成像方程求参考波长的像距与横向放大率，证明 $s'(\lambda)$ 对波长的精确一阶灵敏度；(3) 对小谱宽 $\Delta\lambda$ 求纵向色差 $\Delta s'$ 的一阶近似，并指出红光与蓝光焦点哪个离镜面更远；(4) 求物体处于焦距附近时的像距发散和在薄透镜接触的薄色散极限 $C\to0^+$ 下的结果。
